% chapter: wiki_skin_effect
\chapter{Skin Effect}\label{ch:wiki_skin_effect}

The skin effect is the tendency of alternating electric current to flow primarily near the surface of a conductor. At higher frequencies, the current density decreases exponentially with depth from the surface, concentrating in a thin layer called the skin depth.

				

\section{Skin Depth Formula}

				
\[\delta = \sqrt{\dfrac{2}{\omega\mu\sigma}} = \sqrt{\dfrac{1}{\pi f\mu\sigma}}\]

				

where $\omega = 2\pi f$ is angular frequency, $\mu$ is permeability, and $\sigma$ is conductivity. Current density falls to $1/e \approx 37\%$ of its surface value at depth $\delta$.

\rffig{wiki_skin_effect_1.pdf}{Current density decays exponentially with depth. At depth = δ, J has fallen to 1/e ≈ 37\% of its surface value.}

				

\section{Values for Common Conductors at 100 MHz}

				

\rffig{wiki_skin_effect_2.pdf}{Current density J vs depth d. At one skin depth δ, J falls to 1/e ≈ 37\% of its surface value.}

				
\begin{itemize}

					\item Copper ($\sigma = 5.8\times10^7$ S/m): $\delta \approx 6.6$ µm

					\item Aluminium ($\sigma = 3.8\times10^7$ S/m): $\delta \approx 8.2$ µm

					\item Silver ($\sigma = 6.3\times10^7$ S/m): $\delta \approx 6.3$ µm

				
\end{itemize}

				

\section{Impact on RF Resistance}

				

Because current only flows in a thin annulus near the surface, the effective cross-sectional area of the conductor is reduced, increasing resistance compared to DC. The AC resistance of a round wire of radius $r \gg \delta$ is approximately:

				
\[R_{AC} \approx R_{DC}\cdot\dfrac{r}{2\delta}\]

				

For a 1 mm diameter copper wire at 100 MHz, $R_{AC}/R_{DC} \approx 38$. This is why multi-strand (Litz) wire is used in high-Q inductors at lower RF frequencies — many thin strands, each smaller than the skin depth, provide more effective conductor area.

				

\section{Design Implications}

				
\begin{itemize}

					\item Conductor thickness in PCBs should be at least 3–5 skin depths to minimise resistive loss.

					\item Silver plating adds only marginal benefit over copper above \textasciitilde{}100 MHz (skin depth \textasciitilde{}6 µm) as long as the copper is thicker than a few skin depths.

					\item In MRI coils at 128 MHz (3 T), copper tubing (thick wall) is preferred over solid wire for high-Q resonators.

				
\end{itemize}
